% TOP_PROBLEM: {"id": 1270, "title": "Density of Ulam Numbers", "category_id": 1, "category": {"id": 1, "name": "number_theory", "display_name": "Number Theory", "description": "Properties of integers, prime numbers, Diophantine equations.", "slug": "number-theory", "order_index": 1, "created_at": "2026-07-31T15:26:25.670Z"}, "status": "open"}
% FIRST_POSTED: null
% ENTRY_KIND: statement_only
% Imported from: batch05.tex; UnsolvedMath ID 1270
\documentclass[11pt]{article}
\usepackage[margin=25mm]{geometry}
\usepackage{fontspec}
\setmainfont{Latin Modern Roman}
\usepackage{amsmath,amssymb,mathtools}
\usepackage{unicode-math}
\setmathfont{Latin Modern Math}
\usepackage{xurl,hyperref}
\hypersetup{colorlinks=true,urlcolor=blue}
\setcounter{secnumdepth}{0}
\setlength{\parindent}{0pt}
\setlength{\parskip}{5pt}
\setlength{\emergencystretch}{3em}
\newcommand{\sref}[2]{\hyperlink{source:#1:#2}{[S#2]}}
\newcommand{\cataloguescope}[1]{\paragraph{Recorded target.}#1}
\begin{document}
% UnsolvedMath ID 1270; source catalogue code NT-063
\setcounter{section}{1269}
\section{Density of Ulam Numbers}\label{1270}
{\small\sffamily UnsolvedMath ID 1270 · Number Theory}
\subsection{Definitions and mathematical statement}

\paragraph{Shared notation.}$\mathbb N=\{1,2,\ldots\}$, and $\mathbb Z,\mathbb Q,\mathbb R,\mathbb C$ denote the integers, rationals, reals and complex numbers. Any other conventions are specified locally.


Let $u_1=1$ and $u_2=2$. For each $k\geq3$, let $u_k$ be the least integer strictly greater than $u_{k-1}$ for which exactly one pair of indices $1\leq i<j<k$ satisfies $u_i+u_j=u_k$. Thus order does not distinguish two representations and repetition of the same earlier term is forbidden. Put $U=\{u_k:k\geq1\}$. Its natural density, when it exists, is
\[
d(U)=\lim_{N\to\infty}\frac{\#(U\cap\{1,\ldots,N\})}{N}.
\]
\paragraph{Statement recorded in the supplied source.}
Does the natural-density limit $d(U)$ exist, and is it strictly positive? \sref{1270}{2}

\subsection{Short English statement}
Do the classical Ulam numbers occupy a well-defined positive fraction of the positive integers?
\subsection{Sources}
\begin{description}
\item[\hypertarget{source:1270:1}{S1}] ULAM.AI and the UnsolvedMath Contributors, UnsolvedMath: A Curated Collection of Open Mathematics Problems, record 1270 (NT-063). CC BY 4.0. \url{https://huggingface.co/datasets/ulamai/UnsolvedMath}.
\item[\hypertarget{source:1270:2}{S2}] Stefan Steinerberger, \emph{A hidden signal in the Ulam sequence}, arXiv:1507.00267v6 (2016), abstract and Section 1. \url{https://arxiv.org/pdf/1507.00267}.
\end{description}
\label{end:1270}
\end{document}
